% =========================================================
% 13_applications.tex
% Vectors and Linear Algebra
% =========================================================

% =========================================================
% SECTION DIVIDER
% =========================================================

\section{Applications in Computing and AI}

% =========================================================
% LINEAR ALGEBRA IN COMPUTING
% =========================================================

\begin{frame}{Linear Algebra in Computing}

Linear algebra provides a mathematical language for
representing and transforming structured data.

\vspace{0.4cm}

Common objects include

\begin{itemize}
    \item vectors for features and observations,
    \item matrices for transformations and datasets,
    \item linear systems for solving computational problems,
    \item eigenvectors for discovering important directions.
\end{itemize}

\vspace{0.4cm}

These ideas form part of the mathematical foundation of

\[
\boxed{
\text{Computer Graphics, Machine Learning, AI and Data Science}.
}
\]

\end{frame}


% =========================================================
% DATA REPRESENTATION
% =========================================================

\begin{frame}{Vectors as Data Representation}
\vspace{-0.1cm}
A data instance can be represented as a vector.

For example, a student may be represented by
\vspace{-0.1cm}
\[
\vect{x}
=
\begin{bmatrix}
\text{age}\\
\text{study hours}\\
\text{attendance}\\
\text{previous score}
\end{bmatrix}.
\]

A collection of observations can then be represented
as a matrix
\vspace{-0.1cm}
\[
\boxed{
X=
\begin{bmatrix}
- & \vect{x}_1^T & -\\
- & \vect{x}_2^T & -\\
\vdots & \vdots & \\
- & \vect{x}_m^T & -
\end{bmatrix}.
}
\]
\vspace{-0.1cm}
\[
\boxed{
\text{Data}
\longrightarrow
\text{Vectors}
\longrightarrow
\text{Matrices}.
}
\]

\end{frame}


% =========================================================
% FEATURE VECTORS
% =========================================================

\begin{frame}{Feature Vectors in Machine Learning}

In machine learning, an observation is often represented
by a \textbf{feature vector}.

For example,
\vspace{-0.3cm}
\[
\vect{x}
=
\begin{bmatrix}
x_1\\
x_2\\
\vdots\\
x_n
\end{bmatrix}
\]

may contain $n$ features.

A dataset containing $m$ observations can be represented by
\vspace{-0.1cm}
\[
\boxed{
X\in\mathbb{R}^{m\times n}.
}
\]

Here,
\vspace{-0.1cm}
\[
m=\text{number of observations}
\]

\vspace{-0.1cm}

and
\vspace{-0.1cm}
\[
n=\text{number of features}.
\]

\vspace{-0.1cm}

This representation allows mathematical operations to be
applied to entire datasets.

\end{frame}


% =========================================================
% LINEAR MODELS
% =========================================================

\begin{frame}{Linear Models}

A simple linear model can be written as
\vspace{-0.1cm}
\[
\boxed{
y=\vect{w}^{T}\vect{x}+b.
}
\]

Here,
\vspace{-0.1cm}
\[
\vect{x}
=
\text{input feature vector},
\]
\vspace{-0.1cm}
\[
\vect{w}
=
\text{weight vector},
\]

and
\vspace{-0.1cm}
\[
b
=
\text{bias}.
\]

The dot product
\vspace{-0.1cm}
\[
\vect{w}^{T}\vect{x}
\]

combines the input features using their corresponding
weights.

\end{frame}


% =========================================================
% SYSTEMS IN COMPUTING
% =========================================================

\begin{frame}{Systems of Linear Equations}

Many computational problems can be expressed as

\[
\boxed{
A\vect{x}=\vect{b}.
}
\]

Here,

\[
A=\text{coefficient matrix},
\qquad
\vect{x}=\text{unknown vector},
\qquad
\vect{b}=\text{known vector}.
\]

Applications include

\begin{itemize}
    \item solving networks of equations,
    \item resource allocation,
    \item circuit models,
    \item parameter estimation,
    \item numerical simulation.
\end{itemize}

Gaussian elimination provides a systematic method
for solving such systems.

\end{frame}


% =========================================================
% COMPUTER GRAPHICS
% =========================================================

\begin{frame}{Linear Transformations in Computer Graphics}

Vectors and matrices are fundamental to computer graphics.
\vspace{-0.1cm}
A point or direction can be represented as
\vspace{-0.1cm}
\[
\vect{x}
=
\begin{bmatrix}
x\\y\\z
\end{bmatrix}.
\]

A matrix can transform the vector:
\vspace{-0.1cm}
\[
\boxed{
\vect{x}'=A\vect{x}.
}
\]

Different matrices can represent operations such as

\begin{itemize}
    \item rotation,
    \item scaling,
    \item reflection,
    \item projection.
\end{itemize}

Thus, linear transformations provide a mathematical
framework for manipulating geometric objects.

\end{frame}


% =========================================================
% ROTATION
% =========================================================

\begin{frame}{Rotation as a Matrix Transformation}

A two-dimensional rotation through angle $\theta$ can be
represented by

\[
\boxed{
R=
\begin{bmatrix}
\cos\theta&-\sin\theta\\
\sin\theta&\cos\theta
\end{bmatrix}.
}
\]

For a vector $\vect{x}$,

\[
\boxed{
\vect{x}'=R\vect{x}.
}
\]

The transformation changes the direction of the vector
while preserving its length.

\[
\boxed{
\|\vect{x}'\|=\|\vect{x}\|.
}
\]

Matrix transformations therefore provide a compact way
to describe geometric operations.

\end{frame}


% =========================================================
% IMAGE REPRESENTATION
% =========================================================

\begin{frame}{Matrices and Images}

A grayscale image can be represented as a matrix

\[
I=
\begin{bmatrix}
p_{11}&p_{12}&\cdots&p_{1n}\\
p_{21}&p_{22}&\cdots&p_{2n}\\
\vdots&\vdots&\ddots&\vdots\\
p_{m1}&p_{m2}&\cdots&p_{mn}
\end{bmatrix},
\]

where each entry represents a pixel intensity.

For colour images, multiple matrices can represent
different colour channels.

\vspace{0.4cm}

Therefore,

\[
\boxed{
\text{Image data can be represented using matrices.}
}
\]

Linear algebra can then be used for image transformations,
compression and analysis.

\end{frame}


% =========================================================
% PCA
% =========================================================

\begin{frame}{Principal Component Analysis}

Principal Component Analysis (PCA) uses eigenvectors and
eigenvalues to identify important directions in data.

Given a covariance matrix

\[
C,
\]

the principal directions are obtained from

\[
\boxed{
C\vect{v}=\lambda\vect{v}.
}
\]

The eigenvectors represent principal directions, while
the corresponding eigenvalues indicate the amount of
variation associated with those directions.

PCA can therefore be used for

\[
\boxed{
\text{dimensionality reduction}.
}
\]

\end{frame}


% =========================================================
% PCA INTUITION
% =========================================================

\begin{frame}{PCA: Geometric Intuition}

Suppose data points are distributed mainly along one
direction.

Instead of describing the data using the original
coordinates, PCA finds a new coordinate system aligned
with the dominant directions.

\[
\boxed{
\text{Original coordinates}
\quad\longrightarrow\quad
\text{Principal coordinates}
}
\]

This can reduce the number of dimensions while retaining
important variation in the data.

The new directions are determined by eigenvectors of
the covariance matrix.

\end{frame}


% =========================================================
% DIMENSIONALITY REDUCTION
% =========================================================

\begin{frame}{Dimensionality Reduction}

Suppose
\vspace{-0.1cm}
\[
\vect{x}\in\mathbb{R}^{n}
\]

contains many features.

If the data can be represented approximately using
$k<n$ important directions, we can project it into
\vspace{-0.1cm}
\[
\mathbb{R}^{k}.
\]

Thus,
\vspace{-0.1cm}
\[
\boxed{
n\text{ dimensions}
\longrightarrow
k\text{ dimensions},
\qquad k<n.
}
\]

Benefits can include

\begin{itemize}
    \item reduced computational cost,
    \item lower storage requirements,
    \item visualization of high-dimensional data,
    \item removal of less informative directions.
\end{itemize}

\end{frame}


% =========================================================
% NEURAL NETWORKS
% =========================================================

\begin{frame}{Linear Algebra in Neural Networks}

A basic neural-network layer can be represented as

\[
\boxed{
\vect{z}=W\vect{x}+\vect{b}.
}
\]

Here,

\[
W=\text{weight matrix},
\qquad
\vect{x}=\text{input},
\qquad
\vect{b}=\text{bias}.
\]

A nonlinear activation function is then applied:

\[
\boxed{
\vect{a}=f(\vect{z}).
}
\]

Therefore, neural networks repeatedly combine

\[
\boxed{
\text{matrix multiplication}
+
\text{vector addition}
+
\text{nonlinear transformation}.
}
\]

\end{frame}


% =========================================================
% EMBEDDINGS
% =========================================================

\begin{frame}{Vectors and Embeddings}
\vspace{-0.1cm}
In modern AI systems, objects such as

\begin{itemize}
    \item words,
    \item sentences,
    \item images,
    \item documents,
\end{itemize}

can be represented as vectors called \textbf{embeddings}.

For example,
\vspace{-0.1cm}
\[
\vect{x}\in\mathbb{R}^{d}.
\]

Relationships between embeddings can be studied using
vector operations such as
\vspace{-0.1cm}
\[
\boxed{
\vect{x}\cdot\vect{y}
}
\]

and
\vspace{-0.1cm}
\[
\boxed{
\cos\theta
=
\frac{\vect{x}\cdot\vect{y}}
{\|\vect{x}\|\|\vect{y}\|}.
}
\]

These provide measures of directional similarity.

\end{frame}


% =========================================================
% GRAPH REPRESENTATION
% =========================================================

\begin{frame}{Matrices in Graphs and Networks}

A graph can be represented using matrices.

For example, an adjacency matrix records connections
between vertices:

\[
A_{ij}
=
\begin{cases}
1,&\text{if vertices }i\text{ and }j\text{ are connected},\\
0,&\text{otherwise}.
\end{cases}
\]

Thus,

\[
\boxed{
\text{Graph}
\longrightarrow
\text{Matrix representation}.
}
\]

Matrix operations and eigenvalue methods can then be
used to study structural properties of networks.

\end{frame}


% =========================================================
% RANK IN DATA
% =========================================================

\begin{frame}{Rank and Redundant Information}

Suppose a dataset is represented by

\[
X\in\mathbb{R}^{m\times n}.
\]

If some columns are linear combinations of other columns,
then the features are not all independent.

The rank measures the number of independent directions:

\[
\boxed{
\operatorname{rank}(X)
=
\text{number of independent directions}.
}
\]

A rank smaller than the number of features indicates
linear dependence or redundancy in the representation.

\end{frame}


% =========================================================
% NUMERICAL COMPUTING
% =========================================================

\begin{frame}{Linear Algebra and Numerical Computing}

Many computational algorithms repeatedly perform operations
such as

\[
A\vect{x},
\qquad
A\vect{x}=\vect{b},
\qquad
A^{-1},
\qquad
\det(A).
\]

In practice, efficient numerical algorithms are preferred
over direct symbolic calculations.

Examples include

\begin{itemize}
    \item Gaussian elimination,
    \item LU decomposition,
    \item QR decomposition,
    \item iterative methods.
\end{itemize}

Thus, linear algebra connects mathematical models with
efficient computational procedures.

\end{frame}


% =========================================================
% APPLICATION MAP
% =========================================================

\begin{frame}{Linear Algebra Across AI}

\begin{center}

\begin{tabular}{>{\bfseries}l l}
\toprule
Concept & Example Application\\
\midrule
Vectors & Feature representation\\
Matrices & Datasets and transformations\\
Dot products & Similarity and linear models\\
Linear systems & Numerical modelling\\
Determinants & Invertibility and geometry\\
Eigenvectors & PCA and spectral methods\\
Change of basis & Coordinate transformations\\
Rank & Dimensionality and redundancy\\
Cross products & 3D geometry\\
\bottomrule
\end{tabular}

\end{center}

\vspace{0.4cm}

Linear algebra therefore provides a common mathematical
foundation across many areas of computing and AI.

\end{frame}